% Ad Familiares 3.2 (ad-familiares-03-02)
% Source: https://raw.githubusercontent.com/PerseusDL/canonical-latinLit/master/data/phi0474/phi056/phi0474.phi056.perseus-lat1.xml
% Fetched by scripts/fetch_latin.py

% Dateline (Perseus): Scr. Romae circ. m. Mart. a. 703 (51) .

\ciceroLetterOpener{M. CICERO PROCOS. S. D. APPIO PVLCHRO IMP.}

\ciceroSection{1} Cum et contra voluntatem meam et praeter opinionem accidisset, ut mihi cum imperio in provinciam proficisci necesse esset, in multis et variis molestiis cogitationibusque meis haec una consolatio occurrebat, quod neque tibi amicior, quam ego sum, quisquam posset succedere neque ego ab ullo provinciam accipere, qui mallet eam quam maxime mihi aptam explicatamque tradere. quod si tu quoque eandem de mea voluntate erga te spem habes, ea te profecto numquam fallet. A te maximo opere pro nostra summa coniunctione tuaque singulari humanitate etiam atque etiam quaeso et peto ut, quibuscumque rebus poteris (poteris autem plurimis), prospicias et consulas rationibus meis.

\ciceroSection{2} vides ex senatus consulto provinciam esse habendam. si eam, quod eius facere potueris, quam expeditissimam mihi tradideris, facilior erit mihi quasi decursus mei temporis. quid in eo genere efficere possis, tui consili est; ego te, quod tibi veniet in mentem mea interesse, valde rogo. pluribus verbis ad te scriberem, si aut tua humanitas longiorem orationem exspectaret aut id fieri nostra amicitia pateretur aut res verba desideraret ac non pro se ipsa loqueretur. hoc velim tibi persuadeas, si rationibus meis provisum a te esse intellexero, magnam te ex eo et perpetuam voluptatem esse capturum. vale .
